\documentclass[11pt, letterpaper]{article}

% ── Packages ──
\usepackage[utf8]{inputenc}
\usepackage[T1]{fontenc}
\usepackage{amsmath, amssymb, amsthm, mathtools}
\usepackage{microtype}
% Body + math: Computer Modern (LaTeX/AMS default) via Latin Modern
\usepackage{lmodern}
% Sans-serif (used for section headings): Latin Modern Sans
\renewcommand{\sfdefault}{lmss}
\usepackage{geometry}
\usepackage{hyperref}
\usepackage{enumitem}
\usepackage{booktabs}
\usepackage{multirow}
\usepackage{array}
\usepackage{caption}
\usepackage[svgnames]{xcolor}
\usepackage{titling}
\usepackage{titlesec}

\geometry{margin=1.05in}
\hypersetup{
  colorlinks=true,
  linkcolor=MidnightBlue,
  citecolor=MidnightBlue,
  urlcolor=MidnightBlue,
  pdftitle={A New Lower Bound for the Supremum of Autoconvolutions},
  pdfauthor={Andrei Piterbarg, Jai Bajaj, Derrick Vincent},
  pdfsubject={A computer-assisted lower bound for the autoconvolution constant}
}
\captionsetup{font=small, labelfont=bf}
\setlist{topsep=4pt, itemsep=2pt, parsep=0pt}
\setlength{\parindent}{1.1em}
\setlength{\parskip}{0pt}
\allowdisplaybreaks
\setlength{\droptitle}{-2.8em}
\pretitle{\begin{center}\LARGE}
\posttitle{\par\end{center}\vspace{-0.65em}}
\preauthor{\begin{center}\normalsize}
\postauthor{\par\end{center}\vspace{-0.45em}}
\predate{\begin{center}\small}
\postdate{\par\end{center}\vspace{0.6em}}

% Sections: centered, serif bold (matches the Abstract title font)
\titleformat{\section}[block]
  {\normalfont\Large\bfseries\centering}
  {\thesection.}{0.55em}{}
% Subsections: left-aligned, serif bold (matches the Abstract title font)
\titleformat{\subsection}[hang]
  {\normalfont\large\bfseries}
  {\thesubsection.}{0.5em}{}
\titlespacing*{\section}{0pt}{2.6ex plus 0.8ex minus 0.3ex}{1.4ex}
\titlespacing*{\subsection}{0pt}{1.8ex plus 0.5ex minus 0.2ex}{0.8ex}

% Abstract: match the \section heading style (centered, \Large\bfseries,
% serif) so the "Abstract" title is the same size as "1. Introduction".
% Body keeps the article-class \small + quotation framing.
\renewenvironment{abstract}{%
  \vspace{2.6ex plus 0.8ex minus 0.3ex}%
  \begin{center}{\normalfont\Large\bfseries\abstractname}\end{center}%
  \vspace{1.4ex}%
  \small\quotation
}{%
  \endquotation
}

% ── Theorem environments ──
\newtheoremstyle{paperplain}
  {8pt}{8pt}{\itshape}{}{\bfseries}{.}{0.5em}{}
\newtheoremstyle{paperdef}
  {8pt}{8pt}{\normalfont}{}{\bfseries}{.}{0.5em}{}
\theoremstyle{paperplain}
\newtheorem{theorem}{Theorem}[section]
\newtheorem{lemma}[theorem]{Lemma}
\newtheorem{proposition}[theorem]{Proposition}
\newtheorem{corollary}[theorem]{Corollary}
\newtheorem{claim}[theorem]{Claim}
\theoremstyle{paperdef}
\newtheorem{definition}[theorem]{Definition}
\newtheorem{remark}[theorem]{Remark}
\newtheorem{example}[theorem]{Example}

% ── Macros ──
\newcommand{\R}{\mathbb{R}}
\newcommand{\N}{\mathbb{N}}
\newcommand{\Z}{\mathbb{Z}}
\newcommand{\Q}{\mathbb{Q}}
\newcommand{\Linf}{L^\infty}
\newcommand{\supp}{\operatorname{supp}}
\newcommand{\tv}{\operatorname{tv}}
\newcommand{\rev}{\operatorname{rev}}
\newcommand{\conv}{\operatorname{conv}}
\newcommand{\Karc}{K_{\mathrm{arc}}}
\newcommand{\Kms}{K_{\mathrm{ms}}}
\newcommand{\Mcert}{M_{\mathrm{cert}}}

% ── Hyphenation hints for long compound terms ──
\hyphenation{
  Matolcsi
  Vinuesa
  Bochner-ad-mis-si-ble
  arc-sine
  auto-con-vo-lu-tion
  auto-cor-re-la-tion
  Cloninger
  Stein-er-berg-er
  inter-val
  rig-or-ous
  for-mal-i-za-tion
  multi-scale
  AlphaEvolve
}

% ══════════════════════════════════════════════════════════════════════
\title{\texorpdfstring{%
  A New Lower Bound for the\\
  Supremum of Autoconvolutions
}{A New Lower Bound for the Supremum of Autoconvolutions}}
\author{Andrei Piterbarg, Jai Bajaj and Derrick Vincent}
\date{}

\begin{document}
\maketitle

\begin{abstract}
We establish a lower bound on the autoconvolution constant
$C_{1a}=\inf_f \|f*f\|_{\Linf}/(\int f)^2$, taken over nonnegative
$f\in L^1(\R)$ supported on $(-1/4,1/4)$:
\[
  C_{1a}\;\ge\;1292/1000\;=\;1.292.
\]
This improves on the published lower bound of $1.28$ due to
Cloninger and Steinerberger~\cite{CS17}\footnote{An unpublished
$1.2802$ figure is attributed in Tao's
\texttt{optimizationproblems} repository to Xie (2026, AI-assisted,
unaudited)~\cite{Xie26}; it is not part of the published CS17
result.} and on the rigorous analytic bound of $1.27481$ established
by Matolcsi and Vinuesa~\cite{MV10}.  The argument extends the
Matolcsi--Vinuesa master inequality: the single arcsine kernel is
replaced by a convex combination of three arcsine kernels, and the
trigonometric multiplier $G$ is re-optimized as a $200$-cosine
expansion.  An internal arb-coupled enclosure yields the slightly
sharper $\Mcert\ge 66167/51200\approx 1.29232$, recorded for
completeness but not encoded in the headline Lean theorem (see
Section~\ref{sec:proof}).

The five resulting functionals are evaluated rigorously in
\texttt{flint.arb} interval arithmetic at $256$-bit precision.  The
no-$z_1$ quadratic master inequality, applied to the rounded rational
anchors, yields strict failure at $M=1292/1000$ with margin
$\tau-\Phi(M)\ge 9.6\times 10^{-5}$.  The analytic reduction is
formalized in Lean~4 across $30$ modules ($\approx 15.5$~thousand
lines) rooted at \texttt{lean/Sidon/MultiScale.lean}, with two
machine-checked headline theorems.  A \emph{conditional} headline
\texttt{autoconvolution\_ratio\_ge\_1292\_1000} takes an
\texttt{ExtremiserPrimitives}\,$f$ hypothesis record encoding the
outputs of MO~2009 Lemmas~3.1--3.4 / MV~2010 Lemma~3.1 at
$(f,\Kms)$ --- which the paper discharges by direct citation, exactly
as MV~2010 cited MO~2009 for the single-arcsine instance --- and
reaches Lean's three core axioms together with exactly two numerical
user axioms recording the certifier's outputs at the rational slacks
(the $L^2$ integral of $\Kms$ and the gain ratio $(4/u)\,m_G^{\,2}/S_1$).
A second, \emph{unconditional} headline
\texttt{C1a\_ge\_1292\_unconditional} strengthens this: it takes only
the admissibility hypotheses on $f$ (nonnegativity, support in
$(-1/4,1/4)$, $L^1\cap L^2$, unit mass) and constructs the bundle
internally via \texttt{ExtremiserPrimitives.of\_admissible}, so no
hypothesis record is carried.  Its dependency closure reaches the
same three core axioms together with four numerical user axioms (the
two above, plus the certified lower bound $m_G\ge 998/1000$ and the
active-set lattice positivity $\widetilde{\Kms}(j)>0$ for
$1\le j\le 200$).  The quadratic inversion and the slack-soundness
statements are Lean theorems throughout; see
Section~\ref{subsec:proof-lean}.
\end{abstract}

% ======================================================================
\section{Introduction}\label{sec:intro}
% ======================================================================

\subsection{Overview}\label{subsec:intro-overview}

This subsection is an informal account of what is proved and why it
matters; the precise definitions, the numerical record, and the formal
statement of the result follow in
Sections~\ref{subsec:intro-constant}--\ref{subsec:intro-strategy}.

\paragraph{The question.}
Fix a nonnegative function $f$ on the real line, supported in the short
interval $(-1/4,1/4)$ and normalized so that $\int f=1$; think of $f$
as a probability density.  If $X$ and $Y$ are independent random
variables each with density $f$, then their sum $X+Y$ has density the
autoconvolution $f*f$, itself a probability density, now supported in
$(-1/2,1/2)$.  Since $f*f$ integrates to $1$ over an interval of
length $1$, its peak value $\|f*f\|_{\Linf}$ is at least $1$, with
equality only if $f*f$ were the uniform density on $(-1/2,1/2)$ --- a
shape no admissible $f$ produces.  The \emph{autoconvolution constant}
\[
  C_{1a}
  \;=\;
  \inf_{f}\ \|f*f\|_{\Linf}\Big/\Bigl(\textstyle\int f\Bigr)^{2}
\]
measures \emph{how flat} an autoconvolution can be made: it is the
smallest peak attainable as $f$ ranges over all admissible shapes.
Its exact value is unknown; the problem is to trap it between explicit
numerical bounds, and in particular to push the lower bound upward.

\paragraph{Why it matters.}
This continuous extremal problem is the analytic shadow of a classical
question in additive combinatorics going back to Erd\H{o}s and
Tur\'an~\cite{ET41}: how large can a set of integers be if no integer
has too many representations as a sum of two of its elements (a Sidon,
or $B_2$, condition)?  Martin and O'Bryant~\cite{MO2009} made the link
precise, showing that the asymptotics of such extremal sets are
controlled by this same constant $C_{1a}$.  Each improvement to the
\emph{rigorous} lower bound on $C_{1a}$ sharpens the corresponding
asymptotic results for $B_2^{+}[1]$ and related Sidon variants, which
is why the constant has drawn sustained attention for several decades.

\paragraph{Why lower bounds are hard.}
Exhibiting one function $f$ can only ever yield an \emph{upper} bound
on the infimum $C_{1a}$.  A \emph{lower} bound is a statement of a
different kind: it must certify, simultaneously and for \emph{every}
admissible $f$, that the peak cannot be small.  Matolcsi and
Vinuesa~\cite{MV10} supply exactly such a certificate.  From an
auxiliary pair chosen by us --- a nonnegative ``kernel'' $K$ whose
Fourier coefficients are also nonnegative (a \emph{Bochner-admissible}
kernel) together with a trigonometric ``multiplier'' $G$ --- their
duality produces a single quadratic inequality that every admissible
$f$ must satisfy, and hence a rigorous lower bound on $C_{1a}$.  Any
valid pair $(K,G)$ yields \emph{some} bound; the entire game is to
construct a pair that yields a \emph{good} one.  With a single arcsine
kernel and a $119$-term cosine multiplier, Matolcsi and Vinuesa
obtained $C_{1a}\ge 1.27481$; on the other side, the constant is known
to be at most $\approx 1.502862$~\cite{AlphaEvolve26}.

\paragraph{What is new here.}
We enlarge the class of certificates.  Matolcsi and Vinuesa worked with
a kernel at a \emph{single} scale.  We use that a convex combination of
arcsine kernels at \emph{several} scales is still Bochner-admissible
(Theorem~\ref{thm:admissibility}), and that mixing scales removes a
structural bottleneck.  A single-scale arcsine kernel has ``blind
spots'': frequencies at which its Fourier transform nearly vanishes,
namely the zeros of a Bessel function.  Each near-zero sits in a
denominator of the certificate and inflates a critical quantity
$S_1$, throttling the resulting bound.  Different scales place their
blind spots at different frequencies, so a three-scale blend keeps the
Fourier transform uniformly away from zero: $S_1$ falls from
$\approx 87.4$ to $\le 29.841$, and the bound rises in proportion.
Capturing the gain also requires re-optimizing the multiplier $G$ for
the blended kernel, which we do by convex programming with $200$
cosine terms.  The resulting three-scale construction --- arcsine
half-widths $(138,55,25)/1000$ with weights $(85,10,5)/100$ --- gives
the new lower bound
\[
  C_{1a}\ \ge\ 1292/1000\ =\ 1.292,
\]
(Theorem~\ref{thm:main}).

\paragraph{In what sense this is rigorous.}
Closing the certificate reduces to controlling five real numbers
extracted from $(K,G)$.  These are not merely estimated in floating
point: each is enclosed in a mathematically guaranteed interval using
\texttt{flint.arb} arbitrary-precision interval arithmetic at $256$-bit
precision~\cite{Flint}, and each interval is then rounded
\emph{outward} to an exact rational.  The final quadratic step that
produces $1.292$ is carried out in exact rational arithmetic, so the
conclusion is independent of any floating-point rounding.
The trust base for the numerical step is therefore Lean's logical
core together with the soundness of \texttt{flint.arb}'s interval
arithmetic~\cite{Flint}, a standard assumption in computer-assisted
real-number results (cf.\ Flyspeck for Kepler).

Independently, the entire analytic reduction --- the master
inequality, its inversion, and the admissibility of the three-scale
kernel --- is mechanized in the Lean~4 proof
assistant~\cite{Lean,Mathlib} ($30$ modules, $\approx 15.5$~thousand
lines, no \texttt{sorry}).  The mechanization provides two
machine-checked headline theorems.  The \emph{conditional} headline
rests only on Lean's three logical axioms together with exactly two
explicitly identified numerical inputs: decidable inequalities about
concrete real analytic expressions inside Lean --- the $L^2$ integral
of the three-scale kernel $\Kms$, and the gain ratio
$(4/u)\,m_G^{\,2}/S_1$ formed from the embedded $200$ rational QP
coefficients and the closed-form Bessel-based periodic Fourier
coefficients of $\Kms$ --- that play the role of the ``Mathematica
computed this value'' citations in~\cite{MV10}, here backed by
certified interval arithmetic rather than heuristic numerics.

The standard Fourier-analytic facts of MO~2009 Lemmas~3.1--3.4 /
MV~2010 Lemma~3.1 are applied directly to the admissible
kernel~$\Kms$ (a pdf supported in $[-\delta_1,\delta_1]$ with
nonnegative periodic Fourier coefficients and $\Kms\in L^2$); the
paper discharges those outputs at $(f,\Kms)$ by direct citation,
exactly as MV~2010 cited MO~2009 for the single-arcsine application.
The conditional headline carries those facts as the named fields of
an \texttt{ExtremiserPrimitives}\,$f$ record.  The mechanization then
goes one step further than the paper-level citation: a dedicated
\texttt{Sidon.Constructor} layer ($17$ axiom-free modules) constructs
this record for \emph{any} admissible $f$ --- mechanizing the $L^1$
convolution Fourier identity, $\Kms\in L^2$ via the Young estimate
$L^{4/3}\!\star L^{4/3}\!\to L^2$, a period-$u$ Poisson sampling
lemma (generalizing mathlib's period-$1$/continuous version),
MO~2009 Lemma~2.1 in period-$u$ Parseval form, the period-$1$
Parseval identity for $f\circ f$, and the Cauchy--Schwarz floor.  The
resulting \emph{unconditional} headline
\texttt{C1a\_ge\_1292\_unconditional} takes only the admissibility
hypotheses on $f$ and discharges the bundle internally via
\texttt{ExtremiserPrimitives.of\_admissible}, removing the hypothesis
record entirely (Section~\ref{subsec:proof-lean}).  It thereby
establishes $C_{1a}\ge 1.292$ to at least the rigor of
Matolcsi--Vinuesa~\cite{MV10}: the analytic content is both cited (as
MV did) and now mechanized in Lean, the assembly is carried out in
exact rational arithmetic, and the only inputs taken on trust are
four computer-checked numerical facts of the same \emph{kind} as
MV's Mathematica citations, but verified in \texttt{flint.arb}
interval arithmetic at $256$-bit precision (rigorous, not heuristic)
and corroborated independently in \texttt{mpmath}.

\paragraph{Roadmap.}
Section~\ref{sec:mv} recalls the Matolcsi--Vinuesa master inequality
and its inversion in the form we need.
Section~\ref{sec:kernel} introduces the three-scale arcsine kernel and
proves its admissibility.  Section~\ref{sec:cert} certifies the five
governing functionals as exact rational bounds via interval
arithmetic.  Section~\ref{sec:proof} closes the inequality at the
rational witness $M=1292/1000$, proving the main theorem, and
summarizes the Lean formalization.

\subsection{The Constant \texorpdfstring{$C_{1a}$}{C1a}}\label{subsec:intro-constant}

Let $\mathcal{F}=\{f\in L^1(\R) : f\ge 0,\
\supp(f)\subseteq(-1/4,1/4),\ \int f>0\}$.  (By the
Schinzel--Schmidt reduction~\cite{SS2002}, we may further assume
$f\in L^2(\R)$ without loss of generality when controlling the
ratio~$R(f)$ from below; the period-$1$ Plancherel identity
$\sum_{j\in\Z}|\widehat{f}(j)|^{4}=\|f*f\|_{L^{2}}^{2}$
used in the master inequality's tail bound is consequently
well-defined.)  The \emph{autoconvolution constant} is
\[
  C_{1a}
  \;=\;
  \inf_{f\in\mathcal{F}} R(f),
  \qquad
  R(f) \;:=\; \frac{\|f*f\|_{\Linf}}{\left(\int f\right)^2}.
\]
By homogeneity we may normalize $\int f=1$, in which case
$R(f)=\|f*f\|_{\Linf}\ge 1$.  Heuristically, $R(f)$ measures the
peak height of the self-convolution $f*f$ relative to its $L^1$
mass, so $C_{1a}$ identifies the flattest autoconvolution achievable
in the class $\mathcal{F}$.  This constant is the continuous
analogue of the Sidon-set extremal problem~\cite{ET41,MO2009};
improving its lower bound tightens known asymptotic results on
$B_2^{+}[1]$ sets and related Sidon variants.

\subsection{Known Bounds}\label{subsec:intro-bounds}

Prior lower bounds on $C_{1a}$ are summarized in
Table~\ref{tab:published-bounds}.  Matolcsi--Vinuesa~\cite{MV10}
established the rigorous analytic value $1.27481$ in 2010, and
Cloninger and Steinerberger~\cite{CS17} announced a computational
improvement to $1.28$ via a discrete branch-and-prune search.  An
unpublished improvement to $1.2802$ is attributed in Tao's
\texttt{optimizationproblems} repository~\cite{Xie26} to Xie (2026,
AI-assisted, unaudited).  All subsequent published progress on the
constant has been on the upper-bound side, where a recent AI-assisted
search lowered the ceiling to $1.502862$~\cite{AlphaEvolve26}.  The
present work raises the lower bound to $1.292$
(Theorem~\ref{thm:main}) by refining the kernel class entering the
Matolcsi--Vinuesa dual framework.

\begin{table}[ht]
\centering
\begin{tabular}{lcl}
  \toprule
  \textbf{Bound} & \textbf{Value} & \textbf{Reference} \\
  \midrule
  Upper bound & $\le 1.502862$  & Georgiev et al.~\cite{AlphaEvolve26} (AlphaEvolve) \\
  Upper bound & $\le 1.50972$ & Matolcsi \& Vinuesa~\cite{MV10} \\
  Lower bound & $\ge 1.262$   & Martin \& O'Bryant~\cite{MO2009} \\
  Lower bound & $\ge 1.27481$ & Matolcsi \& Vinuesa~\cite{MV10} \\
  Lower bound & $\ge 1.28$    & Cloninger \& Steinerberger~\cite{CS17} \\
  Lower bound & $\ge 1.2802$  & Xie~\cite{Xie26} (unpublished, AI-assisted, unaudited) \\
  Lower bound & $\ge 1.292$   & \textbf{This work} \\
  \bottomrule
\end{tabular}
\caption[Published bounds on C1a in the normalization used here.]{Published bounds on $C_{1a}$ in the normalization used here.}
\label{tab:published-bounds}
\end{table}

\subsection{Strategy and Main Result}\label{subsec:intro-strategy}

The Matolcsi--Vinuesa~\cite{MV10} dual framework
(Section~\ref{sec:mv}) parametrizes a Bochner-admissible kernel $K$ and
a cosine multiplier $G$, producing a quadratic lower bound on $R(f)$
in five real functionals $(k_1,K_2,S_1,m_G,a)$, of which $(k_1,K_2)$
are extracted from the kernel and $(S_1,m_G,a)$ from the multiplier
(with $a=(4/u)\,m_G^{\,2}/S_1$ a derived gain functional).  MV~2010 takes $K$ to
be the arcsine kernel at half-width $\delta=138/1000$ and $G$ a
$119$-cosine expansion, achieving $R(f)\ge 1.27481$.  The present
construction (i)~replaces the single arcsine by a three-scale convex
combination $\Kms$ at half-widths $(138,55,25)/1000$ with weights
$(85,10,5)/100$; (ii)~re-optimizes $G$ as a $200$-cosine expansion
against $\Kms$; (iii)~certifies the five functionals in
\texttt{flint.arb} interval arithmetic; and (iv)~solves the no-$z_1$
quadratic master inequality with the certified rationals.  The
quadratic inequality $\Phi(M)\ge\tau$ fails strictly at the rational
witness $M=1292/1000$, with margin $9.6\times 10^{-5}$ from the
rational anchors, giving the main result.

\begin{theorem}[Main result]
\label{thm:main}
Let $f:\R\to\R$ be nonnegative with
$\supp(f)\subseteq(-1/4,1/4)$, $\int f>0$, and
$\|f*f\|_{\Linf}<\infty$.  Then
\[
  \frac{\|f*f\|_{\Linf}}{\left(\int f\right)^2}
  \;\ge\;
  \frac{1292}{1000}.
\]
In particular $C_{1a}\ge 1292/1000=1.292$.
\end{theorem}

The proof is assembled in Section~\ref{sec:proof}; all numerical
anchors are reproducible from the accompanying certificate, and a
Lean~4 mechanization is summarized at the end of
Section~\ref{sec:proof}.

% ======================================================================
\section{The Matolcsi--Vinuesa Master Inequality}\label{sec:mv}
% ======================================================================

We summarize the dual framework of~\cite[\S 3]{MV10}; proofs of the
four ingredient lemmas and of the resulting quadratic are recorded
in that reference.  Throughout, $f\in\mathcal{F}$ with $\int f=1$,
$\widehat{h}(\xi)=\int h(x)e^{-2\pi i\xi x}\,dx$, and
$M=R(f)=\|f*f\|_{\Linf}$.  The framework converts the desired lower
bound on $M$ into a quadratic inequality whose coefficients depend on
two auxiliary objects of our choice, a positive-definite kernel $K$
and a trigonometric multiplier $G$, so that any rigorous numerical
control over $(K,G)$ translates into a rigorous lower bound on $M$.

\subsection{Admissible Kernels and Multipliers}\label{subsec:mv-defs}

\begin{definition}[Bochner-admissible kernel]
\label{def:bochner}
Fix $\delta\in(0,1/4]$ and set $u=1/2+\delta$.  A function
$K:\R\to\R$ is \emph{Bochner-admissible at scale $\delta$} if
\begin{enumerate}[label=\textnormal{(K\arabic*)}, leftmargin=2.4em]
\item $K\ge 0$ a.e.\ on $\R$;
\item $K(-x)=K(x)$ for every $x\in\R$;
\item $\supp(K)\subseteq[-\delta,\delta]$ and $\int K=1$;
\item the periodic Fourier coefficients
      $\widetilde{K}(j):=\widehat{K}(j/u)$ satisfy
      $\widetilde{K}(j)\ge 0$ for every $j\in\Z$.
\end{enumerate}
We write $K_2:=\|K\|_{L^2(\R)}^2$.
\end{definition}

\begin{definition}[Admissible cosine multiplier]
\label{def:multiplier}
Fix $u$ and $N\in\N$.  A trigonometric polynomial
$G(x)=\sum_{j=1}^{N}a_j\cos(2\pi jx/u)$ with $a_j\in\R$ is
\emph{admissible} if there is $m_G>0$ with $G(x)\ge m_G$ for every
$x\in[0,1/4]$.  Define
\begin{equation}\label{eq:S1-a}
  S_1 \;=\; \sum_{j=1}^{N}\frac{a_j^{\,2}}{\widetilde{K}(j)},
  \qquad
  a \;=\; \frac{4}{u}\cdot\frac{m_G^{\,2}}{S_1}.
\end{equation}
\end{definition}

\subsection{The Master Inequality and Its Inversion}\label{subsec:mv-master}

\begin{theorem}[Master inequality, $z_1$-free form]
\label{thm:mv-master}
Let $K$ be Bochner-admissible at scale $\delta$, $u=1/2+\delta$,
and $G$ admissible of degree $N$ with constant $m_G$, defining $S_1$
and $a$ by~\eqref{eq:S1-a}.  Then for every $f\in\mathcal{F}$ with
$\int f>0$,
\begin{equation}\label{eq:master}
  M \;+\; 1 \;+\; \sqrt{(M-1)(K_2-1)}
  \;\ge\;
  \frac{2}{u} \;+\; a.
\end{equation}
\end{theorem}

\begin{proof}
Matolcsi--Vinuesa~\cite[Lemma~3.3, Eq.~(10)]{MV10} prove the sharper
form
\begin{equation}\label{eq:mv-sharp}
  M + 1 + 2 z_1^{\,2} k_1
  + \sqrt{(M-1-2z_1^{\,4})(K_2-1-2k_1^{\,2})}
  \;\ge\;
  \frac{2}{u} + a,
\end{equation}
where $k_1=\widehat{K}(1)$ and $z_1=|\widehat{f}(1)|$ are
period-$1$ Fourier coefficients (so that $\widetilde{K}(j)
=\widehat{K}(j/u)$ of Definition~\ref{def:bochner} is the
period-$u$ companion).  The \emph{$z_1$-free} designation reflects
that $z_1$ is a Fourier coefficient of the (unknown) extremal $f$,
so it cannot be controlled analytically; the Cauchy--Schwarz step
below absorbs $z_1$ (together with $k_1$) into the kernel-side
quantity $K_2$, leaving an inequality that depends only on quantities
we can certify.  Their proof uses the arcsine kernel only through
Bochner-admissibility (Definition~\ref{def:bochner}); the same
derivation therefore applies to any admissible $K$, in particular
to the convex combination~$\Kms$.  The radicand of~\eqref{eq:mv-sharp}
is nonnegative by the same Parseval-type bounds as in~\cite[\S 3]{MV10}:
since $f$ is supported in $[-1/4,1/4]$, period-$1$ Plancherel gives
$\sum_{j\in\Z}|\widehat{f}(j)|^{4}=\|f*f\|_{L^{2}}^{2}\le\|f*f\|_{\Linf}=M$,
so $2z_1^{\,4}\le M-1$ for $f\in\mathcal{F}$ with $\int f=1$; similarly
$\sum_{j\in\Z}|\widehat{K}(j)|^{2}=\|K\|_{L^{2}}^{2}=K_2$ gives
$2k_1^{\,2}\le K_2-1$ for any Bochner-admissible $K$ with $\int K=1$.
Inequality~\eqref{eq:master} then follows from~\eqref{eq:mv-sharp}
by the Cauchy--Schwarz estimate
$\sqrt{ac}+\sqrt{bd}\le\sqrt{(a+b)(c+d)}$ applied to the nonnegative
quadruple
$(a,b,c,d)=(2z_1^{\,4},\,M-1-2z_1^{\,4},\,2k_1^{\,2},\,K_2-1-2k_1^{\,2})$.
\end{proof}

\begin{remark}[From MV's $0.5747/\delta$ to the abstract $K_2$, and the $z_1$-absorption]
\label{rem:mv-eq10-translation}
Two translations separate the form~\eqref{eq:mv-sharp} we use from the
literal statement of~\cite[Lemma~3.3, Eq.~(10)]{MV10}, and we record
both explicitly.  First, MV~2010's Eq.~(10) is stated for the
\emph{single-arcsine} specialization: its radicand contains
$\sqrt{0.5747/\delta-1-2k_1^{\,2}}$, where $0.5747/\delta$ is the
Martin--O'Bryant Lemma~3.2 plug-in
$\|\Karc(\delta;\cdot)\|_{L^2}^{2}=\delta^{-1}\int J_0(\pi\xi)^{4}\,d\xi
\le 0.5747/\delta$ for the single arcsine kernel.  We replace
$0.5747/\delta$ throughout by the abstract quantity
$K_2=\|K\|_{L^2(\R)}^{2}$.  This substitution is valid for any
Bochner-admissible $K$ because MV's proof of Eq.~(10)
(\cite[\S 3, proof of Lemma~3.3]{MV10}) uses the kernel only through
(i) $K\in L^2$ with $K_2=\|K\|_{L^2}^{2}$, (ii) $K\ge 0$, and (iii)
the Bochner positivity $\widetilde{K}(j)\ge 0$ --- it never uses the
arcsine shape beyond these three facts, and the value $0.5747/\delta$
enters only as the numerical evaluation of $\|K\|_{L^2}^{2}$ at the
single-arcsine choice.  All three facts hold for $\Kms$
(Definition~\ref{def:bochner}, established in
Section~\ref{sec:kernel}), so the same proof yields~\eqref{eq:mv-sharp}
with $K_2=K_2(\Kms)$.  Second, the passage from the sharp
$z_1$,$k_1$-form~\eqref{eq:mv-sharp} to the $z_1$-free
form~\eqref{eq:master} is the present paper's Cauchy--Schwarz step
$\sqrt{ac}+\sqrt{bd}\le\sqrt{(a+b)(c+d)}$, applied to the nonnegative
quadruple $(2z_1^{\,4},\,M-1-2z_1^{\,4},\,2k_1^{\,2},\,K_2-1-2k_1^{\,2})$
of the proof above; nonnegativity of its entries is exactly the
period-1 Plancherel bounds $2z_1^{\,4}\le M-1$ (for $f*f$) and
$2k_1^{\,2}\le K_2-1$ (for $\Kms$).  In short: the cite to MV Eq.~(10)
supplies the underlying $z_1$,$k_1$ structure; the $0.5747/\delta\to
K_2$ substitution and the $z_1$-absorption are this paper's
translations, each justified above.
\end{remark}

\begin{lemma}[Quadratic inversion]
\label{lem:inversion}
For fixed $K_2\ge 1$, the function
\[
  \Phi(M)\;:=\;M+1+\sqrt{(M-1)(K_2-1)}
\]
is continuous and strictly increasing on $[1,\infty)$ with
$\Phi(1)=2$.  For any $\tau\ge 2$ the equation $\Phi(M)=\tau$ has a
unique solution $M_*\in[1,\infty)$, and every $f\in\mathcal{F}$ with
$\Phi(R(f))\ge\tau$ satisfies $R(f)\ge M_*$.  In particular, with
$\tau=2/u+a$, $C_{1a}\ge M_*$.
\end{lemma}

\begin{proof}
For $M>1$, $\Phi'(M)=1+\sqrt{(K_2-1)/(M-1)}/2>1$, so $\Phi$ is
strictly increasing on $(1,\infty)$ and continuous at $M=1$ from the
right.  Existence, uniqueness, and the conclusion follow from the
intermediate value theorem and Theorem~\ref{thm:mv-master}.
\end{proof}

The MV~2010 single-arcsine optimum yields $M_*\approx 1.2748$.  The
remainder of the paper raises $M_*$ by enlarging the kernel class.

\subsection{Extension to Convex Combinations}\label{subsec:mv-extension}

Matolcsi--Vinuesa~\cite[Lemma~3.1]{MV10} (which itself rests on
Martin--O'Bryant~\cite[Lemmas~3.1--3.4]{MO2009}) was stated for a
single arcsine kernel, but its proof uses $K$ only through the
Bochner-admissibility hypotheses
(K1)--(K4) of Definition~\ref{def:bochner}.  We record the
elementary extension to convex combinations here, since the kernel
$\Kms$ of Section~\ref{sec:kernel} is constructed by precisely such a
combination.  Throughout, $(f\circ f)(x):=\int_{\R}f(t)\,f(x+t)\,dt$
denotes the convolutional autocorrelation in MV's
notation~\cite[\S 3]{MV10}.

\begin{lemma}[MV~Lemma~3.1 extension to convex-combination kernels]
\label{lem:mv-extension}
Let $K=\sum_{i=1}^{n}\lambda_i K_i$ be a finite convex combination
of kernels $K_i$, each Bochner-admissible at some scale
$\delta_i\in(0,1/4]$ in the sense of
Definition~\ref{def:bochner}, with $\lambda_i\ge 0$ and
$\sum_i\lambda_i=1$.  Then $K$ is also Bochner-admissible at scale
$\delta_{\max}:=\max_i\delta_i$, and for every
$f\in\mathcal{F}$ with $\int f=1$ the four atomic identities of
\cite[Lemma~3.1]{MV10} / \cite[Lemmas~3.1--3.4]{MO2009}, together
with the $z_1$-absorption Cauchy--Schwarz step in the proof of
Theorem~\ref{thm:mv-master}, deliver the following bounds at the
pair $(f,K)$:
\begin{enumerate}[label=\textnormal{(\arabic*)}, leftmargin=2.6em]
\item $\displaystyle\int_{\R}(f*f)(x)\,K(x)\,dx \;\le\; \|f*f\|_{\Linf}$.
\item $\displaystyle\int_{\R}(f\circ f)(x)\,K(x)\,dx \;\le\;
       1 + \sqrt{\|f*f\|_{\Linf}-1}\,\sqrt{K_2-1}$.
\item $\displaystyle\frac{2}{u} \;+\; 2u\sum_{j\ne 0}\bigl(\operatorname{Re}\widetilde{f}(j)\bigr)^{2}\widetilde{K}(j)
       \;\le\;
       \int_{\R}\bigl((f*f)(x)+(f\circ f)(x)\bigr)K(x)\,dx$,
       where $u=1/2+\delta_{\max}$,
       $\widetilde{f}(j):=\tfrac1u\int_{-u/2}^{u/2}f(x)e^{-2\pi ijx/u}\,dx$
       is the (period-$u$ averaged) Fourier coefficient of $f$, and
       $\widetilde{K}(j):=\widehat{K}(j/u)$ as in
       Definition~\ref{def:bochner} (see
       Remark~\ref{rem:eq3-inequality}).
\item For every even, real-valued, $u$-periodic trigonometric polynomial
       $G$ with $\widetilde{G}(0)=0$ and $\min_{x\in[0,1/4]}G(x)\ge m_G>0$,
\[
  u^{2}\sum_{j\ne 0}\bigl(\operatorname{Re}\widetilde{f}(j)\bigr)^{2}\widetilde{K}(j)
  \;\ge\;
  \frac{m_G^{\,2}}{\displaystyle\sum_{j:\widetilde{G}(j)\ne 0}\widetilde{G}(j)^{2}/\widetilde{K}(j)}.
\]
\end{enumerate}
\end{lemma}

\begin{proof}
\emph{Preservation of (K1)--(K4) under convex combination.}
Nonnegativity, evenness, and the support inclusion
$\supp K\subseteq\bigcup_i\supp K_i\subseteq[-\delta_{\max},\delta_{\max}]$
are immediate.  Unit mass and Fourier positivity follow from linearity:
$\int K=\sum_i\lambda_i\int K_i=\sum_i\lambda_i=1$, and
$\widetilde{K}(j)=\widehat{K}(j/u)=\sum_i\lambda_i\widetilde{K_i}(j)\ge 0$
since each $\widetilde{K_i}(j)\ge 0$ and $\lambda_i\ge 0$.  Hence
$K$ satisfies (K1)--(K4) at scale $\delta_{\max}$.

\emph{Identities (1)--(4) transport.}
Each identity uses $K$ only through (K1)--(K4):
\begin{itemize}[leftmargin=2em]
\item (1) is the pointwise bound $(f*f)\cdot K\le\|f*f\|_{\Linf}\cdot K$
       integrated using $K\ge 0$ and $\int K=1$ (K1, K3).
\item (2) is MV~2010 Lemma~3.3 / MO~2009 Lemma~3.2 (period-$1$
       Parseval applied to $f\circ f$ paired against $K$, followed by
       a sequence Cauchy--Schwarz, in their sharp $z_1$,$k_1$-form
       Eq.~\eqref{eq:mv-sharp}'s analogue for the $f\circ f$~piece),
       followed by the $z_1$-absorption Cauchy--Schwarz step
       $\sqrt{ac}+\sqrt{bd}\le\sqrt{(a+b)(c+d)}$ of
       Theorem~\ref{thm:mv-master}.  Both steps use $K$ only through
       $\widetilde{K}(j)\ge 0$ (K4) and the finiteness of the energy
       $K_2=\|K\|_{2}^{2}$ of Definition~\ref{def:bochner}.  For a
       convex combination of $L^{2}$ kernels the latter follows from
       Minkowski's inequality; for the specific three-scale $\Kms$ it
       is verified on the Fourier side via Plancherel in
       Lemma~\ref{lem:Kms-L2}.  (We do \emph{not} use
       $K\in L^{\infty}_{\rm loc}$: the arcsine autoconvolution is
       unbounded at the origin.)
\item (3) is the $\ge$-direction of the torus Parseval identity
       \[
         \int_{\R}\bigl((f*f)+(f\circ f)\bigr)K\,dx
         \;=\;\frac{2}{u}+2u\sum_{j\ne 0}\bigl(\operatorname{Re}\widetilde{f}(j)\bigr)^{2}\widetilde{K}(j),
       \]
       which combines the torus Parseval split for $(f*f)+(f\circ f)$
       against $K$, the constant term
       $2u\,(\operatorname{Re}\widetilde{f}(0))^{2}\widetilde{K}(0)=2/u$
       (from $\widetilde{f}(0)=\tfrac1u\int f=\tfrac1u$ and
       $\widetilde{K}(0)=\widehat{K}(0)=\int K=1$ by (K3)), and the
       Fourier expansion of the tail.  We retain
       only the inequality direction; see
       Remark~\ref{rem:eq3-inequality} for why the $\ge$-direction
       alone suffices and follows from $\widetilde{K}(j)\ge 0$ (K4)
       without the full period-$u$ summation.
\item (4) is the discrete weighted Cauchy--Schwarz
       (Sedrakyan / Titu / Engel form)
       $\bigl(\sum a_j b_j\bigr)^{2}\le\bigl(\sum a_j^{2}\widetilde{K}(j)\bigr)\bigl(\sum b_j^{2}/\widetilde{K}(j)\bigr)$
       applied with
       $a_j=\operatorname{Re}\widetilde{f}(j)\sqrt{\widetilde{K}(j)}$ and
       $b_j=\widetilde{G}(j)/\sqrt{\widetilde{K}(j)}$
       (well-defined for $j$ with $\widetilde{K}(j)>0$),
       combined with the inner-product floor
       \begin{equation}\label{eq:inner-product-floor}
         u\cdot\sum_{j\ne 0}\operatorname{Re}\widetilde{f}(j)\,\widetilde{G}(j)\;\ge\; m_G,
       \end{equation}
       established as follows.  By admissibility $f\ge 0$ with
       $\supp f\subseteq(-1/4,1/4)$, and $\int_{\R}f\,dx=1$ under the
       normalization $\int f=1$ (Section~\ref{subsec:intro-constant}).
       Since $f$ is real-valued and even, $\widetilde{f}(j)\in\R$, and
       the period-$u$ Parseval pairing of the $u$-periodic, real,
       even multiplier $G$ against $f$ reads
       \[
         \int_{\R}f(x)\,G(x)\,dx
         \;=\;u\sum_{j\in\Z}\operatorname{Re}\widetilde{f}(j)\,\widetilde{G}(j)
       \]
       (MO~2009 Lemma~2.1, applied to the pairing of $f$, supported in
       $(-1/4,1/4)$, against $G$).  The $j=0$ term is
       $u\operatorname{Re}\widetilde{f}(0)\,\widetilde{G}(0)=(\int f)\,\widetilde{G}(0)$
       (using $\widetilde{f}(0)=\tfrac1u\int f$),
       which vanishes because the cosine multiplier
       $G(x)=\sum_{j=1}^{N}a_j\cos(2\pi jx/u)$ of
       Definition~\ref{def:multiplier} has no constant term, i.e.
       $\widetilde{G}(0)=0$.  Finally $G(x)\ge m_G$ for
       $x\in[-1/4,1/4]\supseteq\supp f$ together with $f\ge 0$ gives
       the pointwise bound $m_G\,f(x)\le f(x)\,G(x)$ for all $x$
       (trivially for $x\notin\supp f$), so that
       \[
         u\sum_{j\ne 0}\operatorname{Re}\widetilde{f}(j)\,\widetilde{G}(j)
         \;=\;\int_{\R}f(x)\,G(x)\,dx
         \;\ge\;m_G\int_{\R}f(x)\,dx
         \;=\;m_G,
       \]
       which is~\eqref{eq:inner-product-floor}.  (The Lean witness is
       \texttt{Sidon.MVLemmas.mv\_inner\_product\_floor}, which
       integrates the pointwise bound $m_G\,f\le f\,G$ and substitutes
       the torus-Parseval identity.)
\end{itemize}
Each step invokes (K1)--(K4) alone; no specific kernel structure is
used.  Hence the identities transport unchanged to any
$K=\sum_i\lambda_i K_i$ satisfying (K1)--(K4).  The formal Lean
witness of this extension is the module \texttt{Sidon.MVLemmas},
where identities (1)--(4) are proven for a \emph{generic} admissible
$K$ (with the four atomic primitives of (2) and (3) carried as
named hypotheses, exactly mirroring the Fourier-analytic ingredients
listed above).
\end{proof}

\begin{remark}[Why item~(3) is invoked only in its $\ge$-direction]
\label{rem:eq3-inequality}
MO~2009 Lemma~3.3 / MV~2010 Lemma~3.3 state item~(3) as an
\emph{equality},
\[
  \int_{\R}\bigl((f*f)+(f\circ f)\bigr)K\,dx
  \;=\;\frac{2}{u}+2u\sum_{j\ne 0}\bigl(\operatorname{Re}\widetilde{f}(j)\bigr)^{2}\widetilde{K}(j),
\]
and the master inequality of Theorem~\ref{thm:mv-master} chains
against this only through the lower bound on the left-hand side, i.e.
the $\ge$-direction printed in Lemma~\ref{lem:mv-extension}(3).  The
equality of course implies the inequality by reflexivity, so the
citation to MO~2009 / MV~2010 remains valid; we record here that the
$\ge$-direction needs strictly less.  Indeed, since
$\widetilde{K}(j)\ge 0$ for every $j$ (K4) and $(\operatorname{Re}\widetilde{f}(j))^{2}\ge 0$,
every term of the tail sum is nonnegative, so for any finite
$J\subseteq\Z\setminus\{0\}$ the truncated identity
\[
  \int_{\R}\bigl((f*f)+(f\circ f)\bigr)K\,dx
  \;\ge\;\frac{2}{u}+2u\sum_{j\in J}\bigl(\operatorname{Re}\widetilde{f}(j)\bigr)^{2}\widetilde{K}(j)
\]
follows from finite-$J$ Parseval and Bochner positivity alone, with
the full sum recovered by monotone convergence; no period-$u$ Poisson
summation over all of $\Z$ is required.  This is precisely the
refactor that lets the Lean formalization
(\texttt{Sidon.MultiScale}, field \texttt{hEq3\_ge}; supported by
\texttt{Sidon.MVLemmas.mv\_eq3\_ge}) discharge item~(3) from
mathlib's $L^{1}\cap L^{2}$ Plancherel and finite-$J$ Parseval
primitives, without a full period-$u$ Parseval API.
\end{remark}

\begin{remark}[Boundary case of MO~2009 Lemma~2.1]\label{rem:mo-boundary}
The period-$u$ Parseval identity invoked in step~(3) above is
MO~2009 Lemma~2.1 applied to the \emph{bilinear} pairing of
$(f*f)+(f\circ f)$ (supported in $(-1/2,1/2)$) against $K=K_{\mathrm{ms}}$
(supported in $[-\delta_{\max},\delta_{\max}]$).  The hypothesis of MO
Lemma~2.1 requires $\alpha_1+\alpha_2\le u$ where $\alpha_1$, $\alpha_2$
are the support half-widths of the two functions in the pairing.  Here
$\alpha_1=1/2$ and $\alpha_2=\delta_{\max}$, so the hypothesis reads
$1/2+\delta_{\max}\le u$, satisfied with equality by our choice
$u=1/2+\delta_{\max}$ --- the same boundary case used by MV~2010 at the
single-arcsine point $\delta=\delta_1=0.138$.  In particular, the
hypothesis is NOT that each of $f*f$, $f\circ f$, $K$ individually fit
in one period $(-u/2,u/2)$; only the bilinear sum-of-supports condition
is needed.

We are explicit that MO~2009 Lemma~2.1 is \emph{stated} with the weak
inequality $\alpha_1+\alpha_2\le u$ and that we invoke it at the
boundary equality $\alpha_1+\alpha_2=u$; the justification of the
equality case below is the present paper's analytic argument, not a
verbatim invocation of MO.  The argument is a passage to $L^2$-class
representatives.  Lemma~2.1 is a statement about $L^2$ equivalence
classes: its bilinear pairing is unchanged when either factor is
modified on a Lebesgue-null set~\cite[\S 2.1]{Folland}.  Now
$K_{\mathrm{ms}}=\sum_i\lambda_i\eta_{\delta_i}*\eta_{\delta_i}$ is
continuous near $\pm\delta_{\max}$ and vanishes there --- near each
endpoint the supports of the two arcsine factors
$\eta_{\delta_i}\subseteq(-\delta_i/2,\delta_i/2)$ meet only at a
single point, so the convolution integral is zero --- and in any case
$\{\pm\delta_{\max}\}$ is Lebesgue-null, so $K_{\mathrm{ms}}$
represents the same $L^2$ class as a function supported in the
\emph{open} interval $(-\delta_{\max},\delta_{\max})$;
likewise $f\in L^2$ may be taken supported in the open
$(-1/4,1/4)$.  For these representatives
$\alpha_1+\alpha_2<u$ holds with strict inequality on the (open)
support sets, and the pairing value is in any case insensitive to the
measure-zero modification.  This is consistent with how MO~2009
themselves treat the boundary: their proof of Lemma~3.3 reads
``\emph{Since the inequality $\tfrac12+\delta\le u$ is satisfied, we
may apply Lemma~2.1}'' at exactly $u=\tfrac12+\delta$, so the equality
case is one already exercised in the source.

The period-1 Parseval identity used in step~(2) is MO~2009
Lemma~2.2, applied to $f\circ f$ and $K$ both supported in
$(-1/2,1/2)$ (the standard $L^2$ Parseval on a single period).
\end{remark}

The master inequality~\eqref{eq:mv-sharp} of
Theorem~\ref{thm:mv-master} is the algebraic combination
(1)+(2)+(3)+(4) of these four identities together with the
Parseval-type bounds $2z_1^{\,4}\le\|f*f\|_{\Linf}-1$ and
$2k_1^{\,2}\le K_2-1$.  Lemma~\ref{lem:mv-extension} therefore
justifies the application of Theorem~\ref{thm:mv-master} to the
three-scale kernel $\Kms$ of Section~\ref{sec:kernel}.

% ======================================================================
\section{The Three-Scale Arcsine Kernel}\label{sec:kernel}
% ======================================================================

\subsection{Single-Scale Arcsine and the Sonine Identity}\label{subsec:kernel-single}

For $\delta\in(0,1/4]$, let
\[
  \eta_\delta(x)
  \;:=\;
  \frac{1}{\delta}\cdot
  \frac{2/\pi}{\sqrt{1-(2x/\delta)^2}}\,
  \mathbf{1}_{|x|<\delta/2}
\]
denote the arcsine density on $(-\delta/2,\delta/2)$ (the
pushforward of the arcsine density on $(-1/2,1/2)$ under
$u\mapsto\delta u$).  By the Mehler--Sonine integral
representation $J_0(z)=(2/\pi)\int_0^1\!\cos(zt)/\sqrt{1-t^2}\,dt$
\cite[\S 10.9.4]{DLMF}, $\widehat{\eta_\delta}(\xi)=J_0(\pi\delta\xi)$.  The \emph{arcsine kernel at scale $\delta$} is the
auto-convolution
\[
  \Karc(\delta;\cdot)
  \;:=\;
  \eta_\delta*\eta_\delta,
\]
which is nonnegative and even, supported on $[-\delta,\delta]$, with
$\int\Karc(\delta;\cdot)=(\int\eta_\delta)^2=1$.  The convolution
theorem then gives
\begin{equation}\label{eq:sonine}
  \widehat{\Karc(\delta;\cdot)}(\xi)
  \;=\;
  \widehat{\eta_\delta}(\xi)^{2}
  \;=\; J_0(\pi\delta\xi)^2
  \;\ge\; 0.
\end{equation}
In particular each $\Karc(\delta;\cdot)$ is Bochner-admissible at
scale $\delta$, with periodic coefficients
$\widetilde{\Karc}(\delta;j)=J_0(\pi j\delta/u)^2$ for $u\ge 2\delta$.

\subsection{The Convex Combination}\label{subsec:kernel-combo}

\begin{definition}[Three-scale kernel]
\label{def:Kms}
Fix
\begin{equation}\label{eq:anchors}
  (\delta_1,\delta_2,\delta_3)
  \;=\;
  \bigl(\tfrac{138}{1000},\tfrac{55}{1000},\tfrac{25}{1000}\bigr),
  \qquad
  (\lambda_1,\lambda_2,\lambda_3)
  \;=\;
  \bigl(\tfrac{85}{100},\tfrac{10}{100},\tfrac{5}{100}\bigr),
\end{equation}
and set $u=1/2+\delta_1=638/1000$.  The
\emph{three-scale arcsine kernel} is
\[
  \Kms(x)
  \;:=\;
  \sum_{i=1}^{3}\lambda_i\,\Karc(\delta_i;x),
  \qquad x\in\R.
\]
\end{definition}

\begin{table}[ht]
\centering
\begin{tabular}{ccccc}
  \toprule
  $i$ & $\delta_i$ & $\lambda_i$ & $u$ & $N$ \\
  \midrule
  $1$ & $138/1000$ & $85/100$ & \multirow{3}{*}{$638/1000$} & \multirow{3}{*}{$200$} \\
  $2$ & $55/1000$  & $10/100$ & & \\
  $3$ & $25/1000$  & $5/100$  & & \\
  \bottomrule
\end{tabular}
\caption[Parameters of the three-scale arcsine kernel and the cosine multiplier.]{Parameters of the three-scale arcsine kernel and the cosine
multiplier.  The period $u$ and the cosine degree $N$ are shared
across the three scales.}
\label{tab:kernel-params}
\end{table}

\subsection{Admissibility}\label{subsec:kernel-admissibility}

\begin{theorem}[Admissibility]
\label{thm:admissibility}
The kernel $\Kms$ is Bochner-admissible at scale
$\delta_1=138/1000<1/4$, with periodic coefficients
\begin{equation}\label{eq:Ktilde}
  \widetilde{\Kms}(j)
  \;=\;
  \sum_{i=1}^{3}\lambda_i\,
       J_0\!\left(\frac{\pi j\delta_i}{u}\right)^{\!2}
  \;\ge\; 0,
  \qquad j\in\Z.
\end{equation}
\end{theorem}

\begin{proof}
We check conditions (K1)--(K4) of Definition~\ref{def:bochner}.
\begin{itemize}[leftmargin=2em]
\item \emph{(K1)~Nonnegativity, (K2)~evenness.}  Each
  $\lambda_i\Karc(\delta_i;\cdot)\ge 0$ and is even; nonnegative sums
  of even nonnegative functions are even and nonnegative.
\item \emph{(K3)~Support and unit mass.}
  $\supp\Karc(\delta_i;\cdot)\subseteq
  [-\delta_i,\delta_i]\subseteq[-\delta_1,\delta_1]$ for $i\ge 1$, so
  $\supp\Kms\subseteq[-\delta_1,\delta_1]$; and
  $\int\Kms=\sum_i\lambda_i\int\Karc(\delta_i)=\sum_i\lambda_i=1$.
\item \emph{(K4)~Fourier positivity.}  Linearity of the Fourier
  transform and~\eqref{eq:sonine} give
  \[
    \widehat{\Kms}(\xi)
    \;=\;
    \sum_i\lambda_i J_0(\pi\delta_i\xi)^2 \;\ge\; 0,
  \]
  hence~\eqref{eq:Ktilde}.\qedhere
\end{itemize}
\end{proof}

Two analytical properties of $\Kms$ used silently in
Sections~\ref{sec:cert}--\ref{sec:proof} merit explicit
statement.  Both are immediate from the convex-combination
structure and the corresponding single-scale facts; we record them as
sub-lemmas to make the dependence on (K1)--(K4) of
Definition~\ref{def:bochner} transparent.

\begin{lemma}[$\Kms\in L^{2}(\R)$]
\label{lem:Kms-L2}
$\|\Kms\|_{L^{2}(\R)}
\;\le\;\sum_{i=1}^{3}\lambda_i\,\|\eta_{\delta_i}*\eta_{\delta_i}\|_{L^{2}(\R)}
\;<\;\infty$.
\end{lemma}

\begin{proof}
By Minkowski's inequality on $L^{2}(\R)$ applied to the finite
combination $\Kms=\sum_i\lambda_i(\eta_{\delta_i}*\eta_{\delta_i})$,
$\|\Kms\|_{2}\le\sum_i\lambda_i\,\|\eta_{\delta_i}*\eta_{\delta_i}\|_{2}$,
so it suffices that each single-scale autoconvolution lies in
$L^{2}(\R)$.  We argue on the Fourier side, where the singularity is
invisible.  By the Mehler--Sonine identity~\eqref{eq:sonine},
$\widehat{\eta_{\delta_i}*\eta_{\delta_i}}(\xi)=J_0(\pi\delta_i\xi)^{2}$,
so by Plancherel and the substitution $t=\delta_i\xi$,
\[
  \|\eta_{\delta_i}*\eta_{\delta_i}\|_{2}^{2}
  \;=\;\int_{\R}J_0(\pi\delta_i\xi)^{4}\,d\xi
  \;=\;\frac{1}{\delta_i}\int_{\R}J_0(\pi t)^{4}\,dt.
\]
The Watson envelope $|J_0(z)|\le\sqrt{2/(\pi z)}$ for $z\ge 1$ gives
$J_0(\pi t)^{4}=O(t^{-2})$ as $|t|\to\infty$, while $J_0(\pi t)^{4}\le 1$
on $[-1,1]$; hence the integral converges and each autoconvolution is
in $L^{2}(\R)$, so the right-hand side is finite.  We emphasize that
$\eta_{\delta_i}*\eta_{\delta_i}$ is \emph{not} bounded: it carries a
logarithmic singularity at the \emph{origin} (where
$(\eta_{\delta_i}*\eta_{\delta_i})(0)=\int\eta_{\delta_i}^{2}=\infty$),
not at the endpoints.  Square-integrability is genuine---$\log^{2}$ is
locally integrable---but is established here via Plancherel, not via
any local boundedness of $\Kms$.
\end{proof}

\begin{lemma}[Bochner positivity of $\widehat{\Kms}$]
\label{lem:Kms-pos}
For every $\xi\in\R$,
\[
  \widehat{\Kms}(\xi)
  \;=\;\sum_{i=1}^{3}\lambda_i\,J_0(\pi\delta_i\xi)^{2}
  \;\ge\; 0.
\]
In particular $\widetilde{\Kms}(j)\ge 0$ for every $j\in\Z$.
\end{lemma}

\begin{proof}
By linearity of the Fourier transform,
$\widehat{\Kms}=\sum_i\lambda_i\widehat{\eta_{\delta_i}*\eta_{\delta_i}}$.
The convolution theorem and the Mehler--Sonine
identity~\eqref{eq:sonine} give
$\widehat{\eta_{\delta_i}*\eta_{\delta_i}}(\xi)
=\widehat{\eta_{\delta_i}}(\xi)^{2}=J_0(\pi\delta_i\xi)^{2}$.
Each summand $\lambda_i J_0(\pi\delta_i\xi)^{2}$ is nonnegative
($\lambda_i\ge 0$, squares are nonnegative), and the sum of
nonnegatives is nonnegative.  Setting $\xi=j/u$ for $j\in\Z$ yields
the periodic statement.
\end{proof}

\noindent
Lemma~\ref{lem:Kms-L2} discharges the finite-$K_2$ hypothesis
underlying identity~(2) of
Lemma~\ref{lem:mv-extension}; Lemma~\ref{lem:Kms-pos} discharges
(K4), which is the only kernel-dependent hypothesis used in (3)
and (4).

The motivation for the multi-scale construction is structural:
in~\eqref{eq:S1-a}, $S_1$ blows up at
frequencies $j$ where the single-scale $J_0(\pi j\delta/u)^2$ is close
to a Bessel zero.  Mixing three scales shifts the zeros of each
summand, so the sum stays bounded away from zero on $1\le j\le 200$
(numerically $\min_j\widetilde{\Kms}(j)\ge 2.08\times 10^{-4}$).  This
drops $S_1$ from $\approx 87.4$ (single-scale) to $\le 29.841$
(Lemma~\ref{lem:S1}), and proportionally raises the gain~$a$.

% ======================================================================
\section{The Rigorous Numerical Certificate}\label{sec:cert}
% ======================================================================

\subsection{The Five Certified Anchors}\label{subsec:cert-anchors}

For the master inequality applied to $(\Kms,G)$ we need rigorous
bounds on five real functionals.  We replace each by a rational
\emph{anchor}: a floor when the functional enters the master
inequality on the side that weakens the bound under understatement
(i.e.\ with a $\ge$ sign), a ceiling on the side that weakens it
under overstatement ($\le$), in both cases obtained by outward
rounding of the certifier's arb interval enclosure.  Anchoring at
exact rationals lets the final inversion (Section~\ref{sec:proof}) be
carried out in exact rational arithmetic, independent of the
floating-point pipeline that produced the enclosures.  The cosine multiplier
$G(x)=\sum_{j=1}^{N}a_j\cos(2\pi jx/u)$ with $N=200$ has rational
coefficients $a_j\in\Q$ with denominator $10^{8}$, obtained by
rounding the solution of the convex QP that minimizes
$\sum_j a_j^{\,2}/\widetilde{\Kms}(j)$ subject to $G\ge 1$ on a
uniform grid of $5001$ points in $[0, 1/4]$ (solved by MOSEK).  The
five lemmas below are arb-interval computations~\cite{Flint}; all
certified outputs are rounded outward to exact rationals.  Only four
of the five anchors ($K_2$, $S_1$, $m_G$, $a$) enter the
$z_1$-free form~\eqref{eq:master} used in Section~\ref{sec:proof};
the kernel-mass moment $k_1$ is recorded because the certifier
internally uses the sharper inequality~\eqref{eq:mv-sharp} (which
sharpens $M_{\rm cert}$ from $1.292$ to $\approx 1.29232$ before
rounding to the rational headline target).

\begin{lemma}[Kernel-mass moment]
\label{lem:k1}
\(
  k_1 \;:=\; \widehat{\Kms}(1)
  \;=\; \sum_{i=1}^{3}\lambda_i J_0(\pi\delta_i)^2
  \;\ge\; 9212/10000.
\)
\end{lemma}

\begin{proof}
At $256$-bit precision, \texttt{arb.bessel\_j(0)} evaluated at exact
rational $\pi\delta_i$ returns balls of radius $<7\times 10^{-77}$;
the certified sum is $k_1\ge 0.92124658$, with slack
$\ge 4.6\times 10^{-5}$ above $9212/10000$.
\end{proof}

\begin{lemma}[Kernel-energy moment]
\label{lem:K2}
\(
  K_2 \;=\; \|\Kms\|_{L^2(\R)}^{2}
  \;\le\; 47897/10000.
\)
\end{lemma}

\begin{proof}
Split
$K_2 = 2\int_0^{T}\widehat{\Kms}(\xi)^2\,d\xi
+ 2\int_{T}^{\infty}\widehat{\Kms}(\xi)^2\,d\xi$ with $T=10^{5}$.
The bulk integral is computed by arb adaptive Gauss--Legendre
quadrature and certified at $4.78882342\ldots$ with arb radius
below $10^{-10}$.  For the tail, the uniform bound
$|J_0(z)|^2\le 2/(\pi z)$~\cite{Watson} holds for $z>0$;
since $\xi\ge T=10^{5}$ in the tail forces $z=\pi\delta_i\xi\ge
\pi(25/1000)(10^{5})\approx 7854$, the bound applies, and applied to
both Bessel factors gives
$J_0(\pi\delta_i\xi)^2 J_0(\pi\delta_j\xi)^2
\le 4/(\pi^4\delta_i\delta_j\xi^2)$, whence
$\widehat{\Kms}(\xi)^2 = \sum_{i,j}\lambda_i\lambda_j
   J_0(\pi\delta_i\xi)^2 J_0(\pi\delta_j\xi)^2
\le 4C^2/(\pi^4\xi^2)$ with
$C=\sum_i\lambda_i/\delta_i\approx 9.978$.  Integrating
$\int_T^\infty \xi^{-2}\,d\xi = 1/T$ and doubling for both half-lines,
$2\int_T^\infty\widehat{\Kms}^2 \le 8C^2/(\pi^4 T) \le 8.19\times 10^{-5}$.
Summing,
$K_2\in[4.78882342,\,4.78890519]\subseteq[0,\,47897/10000]$, slack
$\ge 7.9\times 10^{-4}$.
\end{proof}

\begin{lemma}[Multiplier denominator]
\label{lem:S1}
\(
  S_1 \;=\;
  \sum_{j=1}^{200}\frac{a_j^{\,2}}{\widetilde{\Kms}(j)}
  \;\le\; 29841/1000.
\)
\end{lemma}

\begin{proof}
Each $\widetilde{\Kms}(j)$ is evaluated in arb from~\eqref{eq:Ktilde}
with radius $<10^{-70}$; dividing exact rational squares $a_j^{\,2}$
yields a certified sum $S_1\le 29.840907$, slack
$\ge 9.3\times 10^{-5}$.
\end{proof}

\begin{lemma}[Pointwise lower bound of $G$]
\label{lem:mG}
\(
  m_G \;:=\; \min_{x\in[0,1/4]}G(x) \;\ge\; 998/1000.
\)
\end{lemma}

\begin{proof}
Partition $[0,1/4]$ into $n=32768$ closed cells of half-width
$r=1/(8n)\approx 3.81\times 10^{-6}$.  On each cell centred at $c$,
the second-order Taylor expansion of $G$ yields the rigorous arb
enclosure
\[
  G(\mathrm{cell})\;\subseteq\;
    G(c)\;+\;G'(c)\cdot[-r,r]\;+\;G''(\mathrm{cell})\cdot[-r^{2}/2,\,r^{2}/2],
\]
with $G''$ evaluated as an arb interval over the entire cell from
\[
  G''(x) \;=\; -\sum_{j=1}^{200}a_j\,(2\pi j/u)^{2}\cos(2\pi jx/u).
\]
The minimum of the lower endpoints over all $n$ cells gives the
certified bound $m_G\ge 0.99997987$, slack $\ge 1.9\times 10^{-3}$
above $998/1000$.
\end{proof}

\begin{lemma}[Gain functional]
\label{lem:a}
\(
  a \;=\; (4/u)\cdot m_G^{\,2}/S_1 \;\ge\; 20925/100000.
\)
\end{lemma}

\begin{proof}
Substituting the rational anchors of
Lemmas~\ref{lem:S1} and~\ref{lem:mG} into the definition of $a$
in~\eqref{eq:S1-a},
\[
  a
  \;\ge\;
  \frac{4}{638/1000}\cdot\frac{(998/1000)^{2}}{29841/1000}
  \;=\;
  \frac{4000\cdot 998^{2}}{638\cdot 29841\cdot 10^{3}}
  \;=\;
  \frac{3984016000}{19038558000}
  \;=\; 0.20926\ldots,
\]
hence $a\ge 20925/100000$, slack $\ge 1.0\times 10^{-5}$.  (The
underlying arb interval evaluation returns the tighter enclosure
$a\ge 0.21009214$ when $m_G$ and $S_1$ are coupled rather than
rounded separately; we use only the looser rational floor in what
follows.)
\end{proof}

\begin{table}[ht]
\centering
\begin{tabular}{lcccc}
  \toprule
  Functional & Symbol & Direction & Rational bound & Decimal \\
  \midrule
  Kernel-mass moment   & $k_1$ & $\ge$ & $9212/10000$    & $0.9212$ \\
  Kernel-energy        & $K_2$ & $\le$ & $47897/10000$   & $4.7897$ \\
  Multiplier denominator & $S_1$ & $\le$ & $29841/1000$  & $29.841$ \\
  Multiplier minimum   & $m_G$ & $\ge$ & $998/1000$      & $0.998$  \\
  Gain                 & $a$   & $\ge$ & $20925/100000$  & $0.20925$ \\
  \bottomrule
\end{tabular}
\caption[Certified rational anchors for the three-scale kernel.]{Certified rational anchors for the three-scale kernel
$\Kms$ at parameters~\eqref{eq:anchors} with $u=638/1000$ and the
$200$-cosine multiplier $G$.}
\label{tab:anchors}
\end{table}

\subsection{Rescaling Invariance}\label{subsec:cert-rescaling}

Although the proof of Lemma~\ref{lem:mG}
produces $m_G\ge 998/1000$, the master inequality is invariant
under $G\mapsto G/m_G$: rescaling multiplies both $S_1$ and $m_G^{\,2}$
by $1/m_G^{\,2}$, leaving $a$ unchanged.  The certified rational
bound $m_G\ge 998/1000$ is therefore sufficient for the conclusion,
and the lemma is not sharpened further.

\subsection{Reproducibility}\label{subsec:cert-repro}

All five anchors are produced by the reproducible script and the
associated certificate JSON: each entry is the outward-rounded enclosure
returned by \texttt{flint.arb} at $256$-bit precision; full pipeline
inputs (the rational $a_j$, the bisection schedule, the
quadrature tolerance) and the auditor logs are included in the
accompanying repository.

% ======================================================================
\section{Proof of the Main Theorem}\label{sec:proof}
% ======================================================================

\subsection{The Strict-Failure Witness}\label{subsec:proof-witness}

We close the master inequality using the anchors of
Table~\ref{tab:anchors}.  The strategy is to exhibit a rational
\emph{witness} $M_0=1292/1000$ for which $\Phi(M_0)<\tau$ holds
strictly; combined with the master inequality $\Phi(R(f))\ge\tau$,
this rules out $R(f)\le M_0$ by strict monotonicity of $\Phi$
(Lemma~\ref{lem:inversion}), forcing $R(f)>M_0$.  No analytic
optimization over $M$ is needed: the rational $M_0$ is fixed in
advance and the verification is purely arithmetic.
Set $\tau:=2/u+a$ and recall
$\Phi(M):=M+1+\sqrt{(M-1)(K_2-1)}$ from Lemma~\ref{lem:inversion}.
With the certified rationals from Table~\ref{tab:anchors},
\[
  \tau
  \;\ge\; \frac{2}{638/1000} + \frac{20925}{100000}
  \;=\; 3.134796\ldots + 0.20925
  \;=\; 3.344046\ldots
\]

\begin{proposition}[Strict failure at the rational witness]
\label{prop:fail}
With $K_2\le 47897/10000$ and $\tau=2/u+a\ge 2000/638+20925/10^{5}$,
\begin{equation}\label{eq:strict}
  \Phi(1292/1000) \;\le\; 3.34395
  \;<\; 3.344046\ldots \;\le\; \tau,
\end{equation}
with margin $\tau-\Phi(1292/1000)\;\ge\; 9.6\times 10^{-5}$.
\end{proposition}

\begin{proof}
At $M=1292/1000$, using the certified upper bound $K_2\le 47897/10000$,
\[
  (M-1)(K_2-1) \;\le\; \frac{292}{1000}\cdot\frac{37897}{10000}
  \;=\; \frac{11065924}{10^{7}}
  \;=\; 1.1065924,
\]
and since $\Phi$ is increasing in $K_2$, this upper bound transfers
to $\Phi(M)$.  The rational bound
$(105195/10^{5})^{2}=11065988025/10^{10}>11065924/10^{7}$
gives $\sqrt{(M-1)(K_2-1)}\le 105195/10^{5}=1.05195$, so
\[
  \Phi(1292/1000)
  \;\le\; \frac{1292}{1000} + 1 + \frac{105195}{10^{5}}
  \;=\; \frac{66879}{20000}
  \;=\; 3.34395.
\]
Combining with $\tau\ge 4267003/1276000=3.344046\ldots$ yields
\[
  \tau-\Phi(1292/1000)
  \;\ge\; \frac{4267003}{1276000} - \frac{66879}{20000}
  \;=\; \frac{307}{3190000}
  \;\ge\; 9.6\times 10^{-5}\;>\;0.
\]
The conclusion is independently re-certified by an arb cell-search
bisection at $256$-bit precision, which returns the sharper enclosure
$\Mcert\ge 1.29232$ when the certifier's full arb interval bounds on
all five functionals (rather than the rounded rational anchors of
Table~\ref{tab:anchors}) are used.
\end{proof}

\subsection{Assembly}\label{subsec:proof-assembly}

\begin{theorem}[Restatement of Theorem~\ref{thm:main}]
\label{thm:main-restated}
Let $f\in\mathcal{F}$ with $\|f*f\|_{\Linf}<\infty$.  Then
$R(f)\ge 1292/1000$, and hence $C_{1a}\ge 1292/1000$.
\end{theorem}

\begin{proof}
By homogeneity assume $\int f=1$, so $M:=R(f)=\|f*f\|_{\Linf}\ge 1$.
We assemble in five steps.
\begin{itemize}[leftmargin=2em]
\item[(1)] \emph{Admissibility.}  By Theorem~\ref{thm:admissibility},
$\Kms$ is Bochner-admissible at scale $\delta_1=138/1000$, with
period $u=638/1000$.
\item[(2)] \emph{Anchors.}  By Lemmas~\ref{lem:k1}--\ref{lem:a}, the
multiplier $G$ and the five rational bounds
$k_1\ge 9212/10000$, $K_2\le 47897/10000$, $S_1\le 29841/1000$,
$m_G\ge 998/1000$, and $a\ge 20925/100000$ hold.
\item[(3)] \emph{Master inequality.}  Theorem~\ref{thm:mv-master}
applied to $(\Kms,G)$ yields $\Phi(M)\ge \tau$.  Since $\Phi$ is
increasing in $K_2$, replacing the unknown $K_2$ by the certified
upper bound $47897/10000$ overestimates $\Phi$; the strict failure
$\Phi(M_0)<\tau$ at the witness $M_0=1292/1000$ therefore implies
$\Phi(R(f))>\Phi(M_0)$, which (by Lemma~\ref{lem:inversion}) forces
$R(f)>M_0$.
\item[(4)] \emph{Strict failure.}  By Proposition~\ref{prop:fail},
$\Phi(1292/1000)<\tau$.
\item[(5)] \emph{Conclusion.}  Strict monotonicity of $\Phi$ on
$[1,\infty)$ (Lemma~\ref{lem:inversion}) together with (3) and (4)
forces $M>1292/1000$.  Taking the infimum over $f\in\mathcal{F}$
gives $C_{1a}\ge 1292/1000=1.292$.\qedhere
\end{itemize}
\end{proof}

\subsection{Lean Formalization}\label{subsec:proof-lean}

The analytic chain above is mechanized in Lean~4 across $30$ modules
totalling $\approx 15.5$~thousand lines of code, which build cleanly
under Mathlib~\cite{Mathlib,Lean} with zero \texttt{sorry} tactics.
The formalization is organized in two layers.  The \emph{core} layer
($13$ modules, rooted at \texttt{lean/Sidon/MultiScale.lean}) factors
the analytic reduction into
\texttt{Sidon.\{Defs, Bessel, FourierAux, TorusParseval, MVLemmas,
MasterFromLemmas, BundleDefs, BundleEq1, BundleEq2Schwartz,
BundleEq3Schwartz, BundleEq4, BilinearParseval, MultiScale\}} and
produces the conditional headline theorem
\texttt{autoconvolution\_ratio\_ge\_1292\_1000}.  A \emph{constructor}
layer (\texttt{Sidon.Constructor.\{$\,\ldots\,$\}}, $17$ further
modules, all axiom-free) then assembles the unconditional headline
\texttt{C1a\_ge\_1292\_unconditional}.  We describe the two headlines
and their axiom budgets in turn.

\paragraph{Conditional headline (bundle as hypothesis).}
The theorem \texttt{autoconvolution\_ratio\_ge\_1292\_1000} carries an
\texttt{ExtremiserPrimitives}\,$f$ hypothesis record and reaches
exactly \emph{two} numerical user axioms in its dependency closure
(alongside the three Lean core axioms \texttt{propext},
\texttt{Classical.choice}, \texttt{Quot.sound}):
\begin{itemize}[leftmargin=2em]
\item \texttt{K2\_analytic\_le\_K2UpperQ}: the certified upper bound
$K_2(\Kms)\le\texttt{K2UpperQ}=47897/10000$ from
Lemma~\ref{lem:K2}.  Here \texttt{K2\_analytic} is the concrete real
analytic functional $\int_\R \Kms(x)^2\,dx$ built from the explicit
multi-scale kernel $\Kms$ of
Definition~\ref{def:Kms}.
\item \texttt{gain\_analytic\_ge\_gainLowerQ}: the certified lower bound
$a\ge\texttt{gainLowerQ}=20925/100000$ from Lemma~\ref{lem:a}.  Here
\texttt{gain\_analytic} is the concrete real analytic functional
$(4/u)\cdot m_G^{\,2}/S_1$, with $m_G=\inf_{x\in[0,1/4]}G(x)$ and
$S_1=\sum_{j=1}^{200}a_j^{\,2}/\widetilde{\Kms}(j)$ both formed from
the explicit cosine multiplier $G_{\rm concrete}(x)
=\sum_{j=1}^{200}(\texttt{qpNumerators}[j-1]/10^{8})\cos(2\pi j x/u)$
(the $200$ rational coefficients embedded in \texttt{Sidon.MultiScale}
as the integer list \texttt{qpNumerators}) and the closed-form
Bessel-based periodic coefficients $\widetilde{\Kms}(j)
=\sum_i\lambda_i\,\texttt{besselJ0}(\pi j\delta_i/u)^2$ (using the
power-series Bessel function from \texttt{Sidon.Bessel}).
\end{itemize}
Both axioms bind concrete real expressions in the Lean term
language and are decidable in principle by any sufficient
interval-arithmetic implementation; they play the role of MV~2010's
``Mathematica computed this value'' citations~\cite{MV10}, with the
underlying numerical work performed here by the \texttt{flint.arb}
certifier at $256$-bit precision (in
\texttt{delsarte\_dual/grid\_bound\_alt\_kernel/}).  The three-scale
master inequality at the analytic anchors is itself a Lean
\emph{theorem} (\texttt{MV\_master\_inequality\_from\_MV\_lemmas},
composed from the four MV-Lemma-3.1 atomic conclusions formalized in
\texttt{Sidon.MVLemmas}); slack-monotonicity then lifts it to the
rational anchors used in the strict-failure step.

\paragraph{Unconditional headline (bundle constructed).}
The constructor layer \texttt{Sidon.Constructor} mechanizes the
construction of the \texttt{ExtremiserPrimitives}\,$f$ record from raw
admissibility, so that the bundle is no longer assumed but
\emph{built}.  Its $17$ modules are axiom-free and supply, for any
admissible $f$: the $L^1$ convolution Fourier identity, the
membership $\Kms\in L^2$ via the Young estimate
$L^{4/3}\!\star L^{4/3}\!\to L^2$, a period-$u$ Poisson sampling lemma
(generalizing mathlib's period-$1$/continuous version), MO~2009
Lemma~2.1 in period-$u$ Parseval form, the period-$1$ Parseval
identity for $f\circ f$, and the discrete Cauchy--Schwarz floor of
\eqref{eq:inner-product-floor}.  These are stitched into the bundle by
\texttt{ExtremiserPrimitives.of\_admissible}, and feeding the result
into the conditional headline yields the unconditional theorem
\texttt{C1a\_ge\_1292\_unconditional}.  It takes \emph{only} the
admissibility hypotheses --- $f\in L^1\cap L^2$, $\supp f\subseteq
(-1/4,1/4)$, $f\ge 0$, $\int f=1$ --- and carries no hypothesis record.
Its dependency closure reaches the same three Lean core axioms
together with \emph{four} numerical user axioms: the two above, plus
\begin{itemize}[leftmargin=2em]
\item \texttt{min\_G\_analytic\_ge\_minGLowerQ}: the certified bound
$m_G=\inf_{x\in[0,1/4]}G_{\rm concrete}(x)\ge 998/1000$ from
Lemma~\ref{lem:mG} (a numerical input that the conditional headline
instead absorbs through the bundle's gain field); and
\item \texttt{K\_ms\_fourier\_lattice\_pos\_active}: the certified
active-set positivity $\widetilde{\Kms}(j)>0$ for every integer
$j\in[1,200]$, which licenses the division by $\widetilde{\Kms}(j)$ in
$S_1$ on the Lean side.
\end{itemize}
All four are certifier outputs of the same form --- ``\texttt{flint.arb}
evaluated this functional'' --- decidable in principle by any
sufficient interval-arithmetic implementation.  Because the
constructor layer is axiom-free, the unconditional headline is a
strict strengthening of the conditional one: it reaches the same
mathematical conclusion $C_{1a}\ge 1292/1000$ with the analytic
bundle now \emph{proved} in Lean for every admissible $f$ rather than
\emph{assumed}, at the cost only of the two additional decidable
numerical facts.

\paragraph{Sound definitions of \texorpdfstring{$f\circ f$}{f circ f}
and \texorpdfstring{$\Karc$}{K\_arc}.}
The Lean formalization defines the convolutional autocorrelation as
$(f\circ f)(x):=\int_{\R}f(t)\,f(x+t)\,dt$ (matching the MV
notation~\cite[\S 3]{MV10}; this is \texttt{autocorr} in
\texttt{Sidon.FourierAux}), and the arcsine kernel as the
autoconvolution $\Karc(\delta,\cdot):=\eta_\delta*\eta_\delta$ of the
half-arcsine density $\eta_\delta$ supported on
$(-\delta/2,\delta/2)$.  This matches the analytic definitions used
throughout Sections~\ref{sec:mv}--\ref{sec:kernel} of the present
paper and the Python certifier.

\paragraph{Theorem statements.}
The conditional headline theorem
\texttt{autoconvolution\_ratio\_ge\_1292\_1000} carries a record
\texttt{ExtremiserPrimitives f} whose fields are Lean restatements of
the outputs of MO~2009 Lemmas~3.1--3.4 / MV~2010 Lemma~3.1 at the
pair $(f,\Kms)$; those lemmas apply to~$\Kms$ directly (as a pdf
supported in $[-\delta_1,\delta_1]$ with $\widetilde{\Kms}(j)\ge 0$
and $\Kms\in L^2$), and the paper discharges the bundle fields by
direct citation, exactly as MV~2010 cited MO~2009 at single arcsine.
The unconditional headline \texttt{C1a\_ge\_1292\_unconditional}
removes even the named record: the constructor layer described above
builds it from raw admissibility, so the field-by-field
$L^1\cap L^2$ Plancherel and period-$u$ Parseval content (for which
mathlib~v4.29.1 exposes no single-call API) is supplied by the
purpose-built \texttt{Sidon.Constructor} modules rather than assumed.
The bundle's $m_G$ and $S_G$ fields are reals constrained
by the gain identity $\texttt{gain\_analytic}=2\,m_G^{\,2}/S_G$
(field \texttt{gain\_eq}); the bundle's third inequality field
\texttt{hEq3\_ge} (replacing an earlier equality form \texttt{hEq3})
is Lean-derivable from finite-$J$ Parseval plus Bochner
positivity $\widetilde{\Kms}(j)\ge 0$ alone, via the supporting
theorems \texttt{mv\_eq3\_ge} and \texttt{mv\_eq3\_ge\_of\_eq} in
\texttt{Sidon.MVLemmas} (no period-$u$ Poisson summation required).
A Schwartz-class specialisation that previously bypassed the
\texttt{ExtremiserPrimitives} hypothesis was retired as vacuous:
Paley--Wiener combined with Carlson's theorem prevents any
nontrivial Schwartz $f$ from being compactly supported in
$(-1/4,1/4)$, so the Schwartz headline had no nontrivial witness.

\paragraph{Rigor relative to Matolcsi--Vinuesa.}
Taken together, the two headlines establish $C_{1a}\ge 1.292$ to at
least the standard of rigor of Matolcsi--Vinuesa~\cite{MV10}, and in
several respects beyond it.  The analytic content is discharged both
by direct citation to MO~2009 / MV~2010 applied to $\Kms$ (as MV did
at single arcsine) \emph{and}, in the unconditional headline,
mechanically constructed in Lean for every admissible $f$.  The
numerical content is the same \emph{kind} of computer-checked input
as MV's Mathematica citations, but produced by \texttt{flint.arb}
midpoint--radius interval arithmetic at $256$-bit
precision~\cite{Flint} (rigorous guaranteed enclosures, not heuristic
floating point) and independently corroborated in \texttt{mpmath}.
The final assembly is carried out in exact rational arithmetic with a
strictly positive closing margin: the master inequality's true
threshold is $M_*=1.29203$, so the rational headline target
$1.292$ is a safe round-down, and the strict-failure margin
$\tau-\Phi(1292/1000)\ge 9.6\times 10^{-5}$ of
Proposition~\ref{prop:fail} is verified rationally.  We retain the
honest caveat that four computer-checked numerical axioms remain in
the unconditional headline (two in the conditional one): the result
is computer-assisted in the standard sense (cf.\ Flyspeck for
Kepler~\cite{Flint}), not a pure-kernel hand proof.

Six further statements are Lean \emph{theorems}: the algebraic
inversion \texttt{master\_inequality\_M\_lower} (corresponding to
Proposition~\ref{prop:fail}, proved by case analysis on $M\le 1$ versus
$M>1$ using \texttt{Real.sqrt} monotonicity) and the five
slack-soundness theorems \texttt{K\_two\_upper\_bound},
\texttt{k\_one\_lower\_bound}, \texttt{S\_one\_upper\_bound},
\texttt{min\_G\_lower\_bound}, \texttt{gain\_lower\_bound}, each a
one-line \texttt{norm\_num} verification that the rational slack is on
the correct side of the certifier-reported decimal recorded in
Lemmas~\ref{lem:k1}--\ref{lem:a}.

An explicit finiteness hypothesis $\|f*f\|_{\Linf}<\infty$ appears in
the Lean statement because of the \texttt{ENNReal.toReal} encoding;
this is harmless for the infimum.  The support hypothesis is stated
as $\texttt{Function.support}\,f\subseteq(-1/4,1/4)$, i.e.\
$f$ vanishes pointwise off $(-1/4,1/4)$; for any
$f\in\mathcal{F}$ in the sense of Section~\ref{subsec:intro-constant},
its representative obtained by zero-extending off the essential
support satisfies the pointwise condition and has the same
autoconvolution ratio, so the Lean bound transfers to the
$L^1$-equivalence-class formulation of $C_{1a}$ without loss.

% ======================================================================
\section*{Acknowledgments}
% ======================================================================

We thank M\'at\'e Matolcsi and Carlos Vinuesa, whose dual
framework~\cite{MV10} is the foundation of this argument; and
Alexander Cloninger and Stefan Steinerberger~\cite{CS17}, whose
announced bound of $1.28$ motivated this line of work.  The
formalization rests on the Lean~4 theorem prover~\cite{Lean} and the
Mathlib library~\cite{Mathlib}; the certifier rests on
FLINT/Arb~\cite{Flint}.  We thank the authors and contributors of
each.


% ======================================================================
% References
% ======================================================================

\begin{thebibliography}{ABC10}

\bibitem[GGTW26]{AlphaEvolve26}
Bogdan Georgiev, Javier G\'omez-Serrano, Terence Tao, and Adam~Zsolt
Wagner, \emph{Mathematical exploration and discovery at scale},
preprint, 2026, arXiv:2511.02864,
\url{https://arxiv.org/abs/2511.02864}.

\bibitem[CS17]{CS17}
Alexander Cloninger and Stefan Steinerberger, \emph{On suprema of
autoconvolutions with an application to Sidon sets}, Proc. Amer. Math.
Soc.\ \textbf{145} (2017), no.~8, 3191--3200, arXiv:1403.7988.

\bibitem[ET41]{ET41}
Paul Erd\H{o}s and P\'al Tur\'an, \emph{On a problem of Sidon in additive
number theory, and on some related problems}, J. London Math. Soc.
\textbf{s1-16} (1941), no.~4, 212--215.

\bibitem[Fol99]{Folland}
Gerald B. Folland, \emph{Real Analysis: Modern Techniques and Their
Applications}, 2nd ed., Pure and Applied Mathematics, John Wiley \& Sons,
New York, 1999.

\bibitem[Joh17]{Flint}
Fredrik Johansson, \emph{Arb: efficient arbitrary-precision midpoint-radius
interval arithmetic}, IEEE Trans. Comput. \textbf{66} (2017), no.~8,
1281--1292.  See \url{http://arblib.org/}.

\bibitem[dMU21]{Lean}
Leonardo de Moura and Sebastian Ullrich, \emph{The Lean~4 theorem prover
and programming language}, in: Automated Deduction --- CADE~28, LNCS
\textbf{12699}, Springer, 2021, pp.~625--635.

\bibitem[MO09]{MO2009}
Greg Martin and Kevin O'Bryant, \emph{The supremum of autoconvolutions,
with applications to additive number theory}, Illinois J. Math.
\textbf{53} (2009), no.~1, 219--235, arXiv:0807.5121.

\bibitem[MV10]{MV10}
M\'at\'e Matolcsi and Carlos Vinuesa, \emph{Improved bounds on the supremum
of autoconvolutions}, J. Math. Anal. Appl. \textbf{372} (2010), no.~2,
439--447, arXiv:0907.1379.

\bibitem[The20]{Mathlib}
The mathlib community, \emph{The Lean mathematical library}, Proc.
9th ACM SIGPLAN Int. Conf. Certified Programs and Proofs, 2020,
pp.~367--381.

\bibitem[SS02]{SS2002}
A.~Schinzel and W.~M.~Schmidt,
\emph{Comparison of $L^{1}$- and $L^{\infty}$-norms of squares of polynomials},
Acta Arith. \textbf{104} (2002), no.~3, 283--296.

\bibitem[Wat44]{Watson}
George Neville Watson, \emph{A Treatise on the Theory of Bessel
Functions}, 2nd ed., Cambridge University Press, 1944.

\bibitem[Xie26]{Xie26}
Y.~Xie, \emph{Lower bound entry for $C_{1a}$}, in T.~Tao,
\texttt{optimizationproblems} repository, 2026 (AI-assisted with Grok;
not independently audited).
\url{https://github.com/teorth/optimizationproblems}.

\bibitem[DLMF]{DLMF}
NIST Digital Library of Mathematical Functions,
F.~W.~J. Olver, A.~B. Olde Daalhuis, D.~W. Lozier, B.~I. Schneider,
R.~F. Boisvert, C.~W. Clark, B.~R. Miller, B.~V. Saunders, H.~S. Cohl,
and M.~A. McClain (eds.), \url{https://dlmf.nist.gov/}, Release 1.2.x.

\end{thebibliography}

\end{document}
